% Transcription — fonds Alexandre Grothendieck, shelfmark 133, batch 2 (pages 21-40).
% Established from the facsimile archives/batches/133/batch-02.pdf (a mirror of
% https://grothendieck.umontpellier.fr/133.pdf), per the skill
% transcribe-grothendieck. The batch file carries no cover sheet: PDF page N
% is page 20+N of the folder (the Montpellier footer prints « 21 » ... « 40 »),
% and the archivists' pencil numbers agree wherever legible (« 25 » ... « 33 »,
% « 35 » ... « 40 »).
%
% Pass: Opus 5.5 (claude-opus-5-5), 2026-09-23 — first pass, unchecked against the pages by a human.
% Revised 2026-09-23 (Opus 5.5): header only — pages 21-22 no longer described
% as « for a finite group G (the symmetric groups) »: page 19 has only
% « G groupe », G arbitrary, and the 𝔖_n are the separate factors of
% G × 𝔖_n. The go run re-checked on the facsimile: pages 29-33, five pages,
% all in his hand, none blank. Nothing in the transcription changed.
%
% Eighteen of the twenty pages are transcribed. Absent: page 34, an otherwise
% blank leaf inscribed in his hand « Groupes formels » — a cover, whose words
% become the \section heading before page 35 (references/hand.md §5); page 40,
% a blank leaf, nothing in his hand. No offprint (« tiré à part ») falls in
% this batch, so the folder-22 rule for printed matter is not needed here. He
% did not paginate these leaves. All \section headings but the last are
% bracketed: they are ours, not his; the folder title itself, « [Groupes de
% Witt et formes quadratiques, groupes formels] », is the archivists'.
%
% What the batch is. Five runs, undated, in felt pen (21-22, 29-33), blue ink
% (26-28, 35-39) and black ink (23-25).
%   21-22  Continuation of a run begun in batch 1: for a group G (page 19
%          writes only « G groupe », with no finiteness condition; the
%          symmetric groups enter as the separate factors 𝔖_n of
%          G × 𝔖_n), T^n(E) = cl_n(Hom(I_n, E)) in k_n(G) and the
%          total T^* : k_0(G) -> 1 + k_*(G)^+, multiplicative on sums;
%          Hom(I_n,E) decomposed along the quotients Q of I_n of type
%          α = (α_1..α_n), Σ iα_i = n; a boxed formula expressing T^n through
%          the τ^i by transfer maps φ_{n,i}, cancelled by two large crosses;
%          page 22 is a scratch sheet (the count n!/Π α_i!(i!)^{α_i} of the
%          quotients of type α, a table in τ^n(e_i)). Page 21 opens on « On
%          définit aussi », so its notation (k_n, cl_n, τ^i, I_n) was set up
%          on the previous leaves.
%   23-26  The category ℰ of finite sets and injections; functors ℰ -> (Ensf):
%          representables h_I = Mon(I, -), the functors H_I = E^I, « foncteurs
%          spéciaux » (finite sums of h_I, or direct summands of H_I),
%          « foncteurs distingués » (finite sums of Mon(I,E)/H, H ⊂ 𝔖_I; in a
%          blue addition, ∐_n Mon(I_n,E) ∧^{𝔖_n} X_n), their stability under
%          limits and composition, « équivalences » (isomorphisms from some
%          cardinality on); page 26: Ω = End(ℰ,ℰ) ⊂ Hom(ℰ,Ens), a non-full
%          subcategory, with its operations +, ×, ∘.
%   27-28  On a ringed topos X, K(X,*) = Σ_n K(B_{𝔖_n X}); τ^*(ξ) and the
%          relation τ^n(ξ) = Σ_α Σ_α(λ^1ξ..λ^nξ) ρ_1^{α_1}..ρ_n^{α_n}, with
%          ρ_i the unit of R(𝔖_i) and Σ_α the isobaric polynomial expressing
%          a monomial symmetric function in the elementary ones; corollary:
%          λ^n and the Adams operation as coefficients; question on
%          R_k(1,*) = Σ R_k(𝔖_i) being the polynomial ring on the ρ_i.
%   29-33  The game of go, in his hand: scoring on a ν×ν goban (stones on
%          board, prisoners, territories; D = 2(P+T) - ν^2 - ε, ε = 0 or 1 by
%          who plays last); the rules of capture (« Ko », « remplissage d'un
%          trou »), the end of the obligation to play, the « inattaquable »
%          two-liberty grouping; end-game count by groups with two eyes.
%          Page 33 is a bifolium whose left half is blank.
%   34     Cover « Groupes formels » (absent, see above).
%   35-39  Formal groups: (I) an extension of a p-divisible group by a smooth
%          (or unipotent) one is trivial; a-toric p-divisible subgroups are
%          central; (II) extensions of a smooth formal group by a p-divisible
%          one; DG ∩ Z(G) quasi-unipotent; « Solution »: the invariants
%          L_1 = Z(G), D_1 = D(G), L_0 = G/Z(G), D_0 = G/D(G),
%          Π_1 = L_1 ∩ D_1, Π_0 = G/L_1D_1, with obstruction in H^3(Π_0,Π_1)
%          and indeterminacy in H^2(Π_0,Π_1); page 36 continues with the
%          general crossed situation 1 -> Π_1 -> G_1 -> G_0 -> Π_0 -> 1 and
%          the map Cent(D_0) ∩ Ker θ -> Hom(Π_0,Π_1), whence « Cent D_0 ∩ Ker θ
%          est quasi-unipotent ! ». Pages 37-39 are scratch sheets (boxed
%          unipotence statements, central extensions 1 -> ℨ -> D -> Ĥ -> 1,
%          cocycles φ, ψ for extensions of Ĝ_a by Ĝ_a). Page 35 ends « TSVP »
%          and page 36 does continue it.
%
% Notation fixed for the batch: 𝔖_n as \mathfrak{S}_n; his gothic Z (centre
% of a formal group) as \mathfrak{Z}; ℰ as \mathcal{E}. Words he underlines
% are set in \emph.
\documentclass[11pt,a4paper]{article}
% Le préambule commun à toutes les transcriptions du fonds Grothendieck.
%
% Il tient en une idée : **l'incertitude se compose, elle ne se résout pas.**
% Une transcription qui lisse un mot douteux a détruit la seule chose pour
% laquelle on la faisait — un lecteur doit pouvoir savoir, à chaque endroit,
% ce qui était sur le feuillet et ce qui a été ajouté. Les six macros qui
% suivent sont donc l'appareil critique, et non de la décoration.
%
% Les mêmes commandes sont reconnues par `scripts/render.mjs`, qui produit la
% vue de lecture du site. Une macro ajoutée ici doit l'être aussi là-bas,
% faute de quoi le rendu échouera bruyamment — ce qui est voulu : un rendu
% silencieusement faux est pire qu'un rendu absent.

\ProvidesPackage{grothendieck}[2026/08/08 Transcriptions du fonds Grothendieck]

\usepackage{amsmath,amssymb,amsthm}
\usepackage{mathrsfs}
\usepackage{stmaryrd}
% Les diagrammes commutatifs sont partout dans ce fonds : c'est la langue
% même de la géométrie algébrique de Grothendieck, et une transcription qui
% les rendrait en prose ne transcrirait rien.
\usepackage{tikz-cd}
% Français par défaut — c'est la langue des feuillets — mais l'anglais est
% chargé aussi : la traduction et le résumé sont des documents anglais, et la
% typographie française (espaces avant les deux-points) y serait une faute.
% Ces documents basculent par \selectlanguage{english} en tête de corps.
\usepackage[english,french]{babel}
\usepackage{fontspec}
\usepackage{geometry}
\geometry{a4paper,margin=25mm,marginparwidth=28mm}
\usepackage{marginnote}
\usepackage{xcolor}
% ulem, not soul. `soul`'s \st and \ul are built on a character-by-character
% scan that XeLaTeX's font machinery breaks: the document fails with an
% undefined control sequence rather than degrading. ulem does the same two jobs
% and is Unicode-engine safe. `normalem` stops it hijacking \emph, which the
% transcriptions use for Grothendieck's own emphasis.
\usepackage[normalem]{ulem}
\usepackage{hyperref}
\hypersetup{colorlinks=true,linkcolor=black,urlcolor=black,pdfborder={0 0 0}}

% --- Métadonnées du lot -----------------------------------------------------
% Reprises telles quelles par le script de rendu, qui les affiche en tête de
% la vue de lecture. Elles fixent ce que le fichier prétend couvrir : sans
% elles, une transcription flotte sans point d'attache dans un fonds de seize
% mille feuillets.

\newcommand{\folder}[1]{\gdef\grothFolder{#1}}
\newcommand{\batch}[1]{\gdef\grothBatch{#1}}
\newcommand{\leaves}[2]{\gdef\grothFirst{#1}\gdef\grothLast{#2}}
\newcommand{\foldertitle}[1]{\gdef\grothFoldertitle{#1}}
\gdef\grothFolder{?}\gdef\grothBatch{?}\gdef\grothFirst{?}\gdef\grothLast{?}\gdef\grothFoldertitle{}

% --- Le résumé --------------------------------------------------------------
%
% Il ouvre la lecture modernisée, sous le titre « Résumé », et il est composé
% en petit. Ce n'est pas un ornement : c'est le seul passage du document qui
% ne soit pas des mathématiques, et un lecteur doit pouvoir le reconnaître
% avant de l'avoir lu — puis le sauter s'il connaît déjà le sujet, ou s'y
% arrêter s'il ne le connaît pas. Le filet qui le suit marque la reprise du
% corps.

\newenvironment{resume}
  {\small\setlength{\parskip}{0.45em}}
  {\par\medskip\hrule\medskip}

% --- Mots-clés ---------------------------------------------------------------
%
% En anglais, en clôture du résumé de la lecture modernisée : le vocabulaire
% moderne sous lequel chercher ce que les feuillets construisent. Le manifeste
% du site (`npm run manifest`) les extrait pour étiqueter la cote entière —
% cette ligne est la source unique des tags, il n'y a pas de fichier à côté.
\newcommand{\keywords}[1]{\par\smallskip\noindent
  {\footnotesize\textsc{Keywords} --- \textit{#1}}\par}

% --- L'appareil critique ----------------------------------------------------

% Le numéro de feuillet, en marge. C'est l'ancre qui permet de régler un
% désaccord sur une formule en regardant *un* feuillet plutôt qu'une chemise
% de six cents. Le site s'en sert aussi pour tourner les pages du fac-similé
% quand on fait défiler la transcription.
\newcommand{\leaf}[1]{%
  \marginnote{\footnotesize\color{gray!70}\textsf{#1}}%
  \label{leaf:#1}\ignorespaces}

% Illisible : on ne devine pas. Un mot inventé qui se lit comme les autres est
% le pire résultat possible.
\newcommand{\ill}{\textcolor{red!60!black}{[\,\ldots\,]}}

% Lecture douteuse : le mot est donné, et signalé comme incertain.
\newcommand{\uncertain}[1]{\uline{#1}}

% Ajout de l'éditeur — une lettre manquante, un mot sous-entendu.
\newcommand{\add}[1]{\textcolor{blue!50!black}{[#1]}}

% Biffé par Grothendieck lui-même. Ce qu'il a rayé fait partie du document :
% une rature dit souvent où la pensée a changé de direction.
\newcommand{\struck}[1]{\sout{#1}}

% Note du transcripteur, sur l'état matériel du feuillet.
\newcommand{\note}[1]{\par\smallskip{\small\color{gray!60!black}\textit{[#1]}}\par\smallskip}

% Note marginale de Grothendieck, distincte du corps du texte.
\newcommand{\marginal}[1]{\par\smallskip\noindent\hspace{1em}%
  {\small\color{gray!60!black}#1}\par\smallskip}

% --- Titre ------------------------------------------------------------------

\AtBeginDocument{%
  \begin{center}
    {\large\bfseries Fonds Alexandre Grothendieck --- cote n\textsuperscript{o}~\grothFolder}\\[2pt]
    {\itshape\grothFoldertitle}\\[4pt]
    {\small feuillets \grothFirst--\grothLast\ (lot \grothBatch)}\\[2pt]
    {\footnotesize\color{gray!60!black}Transcription établie d'après le fac-similé numérisé
     par l'Université de Montpellier.\\
     Les lectures incertaines sont soulignées, les illisibles notées [\,\ldots\,],
     les ajouts entre crochets.}
  \end{center}
  \vspace{1em}}

\endinput


\folder{133}
\batch{2}
\pages{21}{40}
\dating{1974-[à partir de 1975]}
\watermark{Édition de démonstration}
\foldertitle{[Groupes de Witt et formes quadratiques, groupes formels] : notes manuscrites (s.d.), tirés à part (1974).}

\begin{document}

\section{[Les opérations $T^n$ et $\tau^i$, suite]}
\note{titre de l'éditeur, entre crochets : le feuillet 21 n'en porte aucun et
commence au milieu d'un développement (« On définit aussi »), dont les
notations $k_n(G)$, $\mathrm{cl}_n$, $\tau^i$, $I_n$ ont été posées sur les
feuillets précédents, au lot 1.}

\page{21}
On définit aussi
\[
  T^n(E) = \mathrm{cl}_n\bigl(\mathrm{Hom}(I_n, E)\bigr) \in k_n(G)
  \qquad (E \in \uncertain{\mathrm{Ob}\,\mathrm{Ens}(G)})
\]
\[
  T^*(E) = \sum_n T^n(E) \in 1 + \widehat{k_*(G)}^{+} ;
\]
on a donc \uncertain{évidemment} $T^*(E \sqcup F) = T^*(E)\,T^*(F)$, d'où
\[
  T^* : k_0(G) \to 1 + \widehat{k_*(G)}^{+}
\]
défini de façon unique par $T^*(\mathrm{cl}_0(E)) = T^*(E)$. On va donner une
relation entre $T^*$ et $\tau^*$, permettant de calculer $T^*$ en termes de
$\tau^*$.
\[
  \mathrm{Hom}(I_n, E) \simeq \coprod_{Q \text{ quotient de } I_n}
  \mathrm{Mon}(Q, E) \simeq
  \coprod_{\substack{\alpha = (\alpha_1, \ldots, \alpha_n) \in \mathbf{N}^n \\
  \sum i\alpha_i = n}}
  \struck{\ill{}}
  \underbrace{\coprod_{Q \text{ quotient de } I_n \text{ de type } \alpha}
  \mathrm{Mon}(Q, E)}_{\varphi_{\alpha}(\tau^{i}(E))}
\]
\note{l'indice de $\varphi$ sous l'accolade est surchargé ; « $\alpha$ » est
une lecture douteuse.}

\uncertain{d'où}
\[
  T^n(E) = \sum_{1 \leqslant i \leqslant n} \varphi_{n,i}\bigl(\tau^i(E)\bigr),
\]
où $\varphi_{n,i} : k_i(G) \to k_n(G)$ défini par
$\varphi_{n,i}(\mathrm{cl}_i(E)) = \mathrm{cl}_n(E \times^{\mathfrak{S}_i}
\mathfrak{S}_n)$ ;
\[
  T^*(E) = \sum_{n \geqslant 0} T^n(E)
  = 1 + \sum_{1 \leqslant i \leqslant n} \varphi_{n,i}\,\tau^i(E)
  = 1 + \sum_{i \geqslant 1} \varphi_{*i}\,\tau^i(E)
\]
\note{tout ce bloc, depuis « d'où », est encadré et annulé par deux grandes
croix ; on le transcrit sans le biffer ligne à ligne. Il se prolonge en haut
du feuillet 22.}

\page{22}
où on pose $\varphi_{*i} : k_i(G) \to \widehat{k_*(G)}^{\geqslant i}$,
$\varphi_{*i}(x) = \sum_{n \geqslant i} \varphi_{ni}(x)$
\note{suite du bloc annulé du feuillet 21, encadré et biffé de même ; une
lettre en exposant sur le premier $\varphi$ ne se lit pas.}

\note{Le reste du feuillet est une page de calculs épars, sans phrase suivie ;
on en donne les éléments dans l'ordre de la page.}

$I_n$ ; \quad $1\cdot\alpha_1 + 2\alpha_2 + \cdots + n\alpha_n = n$, \quad
$\alpha = (\alpha_1, \ldots, \alpha_n)$, $\alpha_i \in \mathbf{N}$
($1 \leqslant i \leqslant n$) ; $E$ un \uncertain{$\mathfrak{S}_i$}-ens.\ fini.

\[
  \begin{array}{ll}
    \alpha_1 \text{ classes de 1 él.} & \alpha_1! \\
    \alpha_2 \text{ — 2 él.} & \alpha_2!\,2^{\alpha_2} \\
    \alpha_i \text{ — } i \text{ él.} & \alpha_i!\,(i!)^{\alpha_i} \\
    \alpha_n \text{ — } n \text{ él.} & \alpha_n!\,(n!)^{\alpha_n}
  \end{array}
\]

$\sum \alpha_i = m$ ; \qquad
$\dfrac{n!}{\prod_{1 \leqslant i \leqslant n} \alpha_i!\,(i!)^{\alpha_i}}$ ;
\qquad $\sum_{Q \text{ quotient de type } \alpha \text{ de } I} E^{\widetilde{Q}}$ ;

\[
  \sum_{\alpha = (\alpha_1, \alpha_2, \ldots) \in \mathbf{N}^{(\mathbf{N}^*)}}
  \varphi_{\nu(\alpha), \mu(\alpha)}\bigl(\tau^{\mu(\alpha)}(x)\bigr)
\]
\note{sous le signe somme, une première indexation
$\alpha_1, \alpha_2, \ldots, \alpha_n$ est biffée ; le premier indice de
$\varphi$ est surchargé et un indice sous $\tau$ est biffé.}

\struck{\ill{}${}^{\alpha_1 \ldots \alpha_i \ldots}(E)$} ; \quad
$\mathbf{Z}[\tau^n(e_i),\ 1 \leqslant i \leqslant d] / \struck{\ill{}}\,
\mathbf{Z}$ ; \quad \struck{$\tau^1(e_1)$} ; \quad un tableau à trois colonnes
$e_i \mid \tau^2(e_i) \mid \tau^3(e_i)$, laissé vide.

\section{[Foncteurs de la catégorie des ensembles finis et injections]}

\page{23}
Soit $\mathcal{E}$ la catégorie des ensembles finis, avec comme morphismes les
applications \emph{injectives}. On considère des foncteurs
$\mathcal{E} \xrightarrow{F} (\mathrm{Ens})$ tels que
$\forall E \in \mathrm{Ob}\,\mathcal{E}$, $\mathrm{card}(F(E))$ soit fini,
i.e.\ des foncteurs $\mathcal{E} \to (\mathrm{Ensf})$. Ils forment une
catégorie
$\overbrace{\underline{\mathrm{Hom}}(\mathcal{E}, (\mathrm{Ensf}))}^{\mathcal{H}}$,
qui est un \uncertain{prétopos}.

Parmi ces objets, on en distingue deux sortes particulières :

a) Les foncteurs \emph{représentables} $\mathcal{E} \to (\mathrm{Ens})$, de la
forme
\[
  h_I : E \mapsto \mathrm{Hom}_{\mathcal{E}}(I, E) = \mathrm{Mon}(I, E)
\]
où $I$ est un ensemble fini. Le foncteur
\[
  I \mapsto h_I \qquad \mathcal{E}^{\circ} \to
  \underline{\mathrm{Hom}}(\mathcal{E}, (\mathrm{Ensf}))
\]
est pleinement fidèle, et identifie $\mathcal{E}^{\circ}$ à une
sous-catégorie pleine de
$\underline{\mathrm{Hom}}(\mathcal{E}, (\mathrm{Ensf}))$.

b) Les foncteurs de la forme
\[
  H_I : E \mapsto \mathrm{Hom}_{(\mathrm{Ens})}(I, E) = E^I = H(I, E)
\]
(restrictions à $\mathcal{E} \subset (\mathrm{Ensf})$ des foncteurs
correspondants sur $(\mathrm{Ensf})$).

\emph{Prop.} Les $h_I \in \mathrm{Ob}\,\underline{\mathrm{Hom}}(\mathcal{E},
(\mathrm{Ensf}))$ sont des objets connexes (non vides) de $\mathcal{H}$. À
isomorphisme

\page{24}
près, ce sont les composantes connexes des foncteurs de la forme $H_J$.

\emph{Foncteurs spéciaux} $\mathcal{E} \to (\mathrm{Ensf})$ : ce sont les
foncteurs qui sont \struck{sommes} isom.\ à des sommes directes finies de
foncteurs de la forme $h_I$ — ou encore qui sont des sommandes directs (pour
$I$ de cardinal assez grand) de foncteurs de la forme $H_I$.

\emph{Stabilités.} Toute somme finie, \struck{\ill{}} produit fini (plus gén.\
lim.\ finies) de foncteurs spéciaux est spécial. Tout foncteur spécial
$\mathcal{E} \to (\mathrm{Ensf})$ prend ses valeurs dans la sous-catégorie
$\mathcal{E}$ de $(\mathrm{Ensf})$, i.e.\ définit un \struck{dans}
$\in \underline{\mathrm{Hom}}(\mathcal{E}, \mathcal{E})$ (et on identifiera
par abus de langage les foncteurs spéciaux $\mathcal{E} \to (\mathrm{Ensf})$
avec \uncertain{des foncteurs} $\mathcal{E} \to \mathcal{E}$). Avec cette
convention, le composé de deux foncteurs spéciaux est spécial.

\page{25}
\emph{Foncteurs distingués} : sommes directes finies de quotient d'un
foncteur \struck{spécial par} $h_I$ par un groupe d'automorphismes, i.e.\
foncteurs isomorphes à des sommes finies de foncteurs de la forme
\[
  E \mapsto \mathrm{Mon}(I, E) / H
\]
où $H \subset \mathfrak{S}_I$. Ce sont aussi les sommes finies de composantes
connexes de foncteurs de la forme
\[
  H(I, E)/H = \mathrm{Hom}_{\mathrm{Ens}}(I, E)/H .
\]
Ces foncteurs ont les mêmes propriétés de permanence que celles vues pour les
foncteurs spéciaux.
\note{Les deux insertions « sommes directes finies de » et « sommes finies
de » sont de sa main, en interligne, rattachées au texte par une accolade.}

\marginal{Tout foncteur distingué \struck{spécial} est isom.\ à un foncteur de
la forme
\[
  E \mapsto \coprod_{n \geqslant 0} \mathrm{Mon}(I_n, E)
  \wedge^{\mathfrak{S}_n} X_n
\]
où $\forall n$, $X_n$ est un $\mathfrak{S}_n$-ens.\ fini
(\uncertain{i.e.}\ \ill{}).}
\note{ajout à l'encre bleue, en regard de la phrase précédente, séparé d'elle
par un trait oblique ; la fin de la parenthèse, en bout de ligne, est
recouverte par une insertion à l'encre noire.}

Un hom.\ $u : F \to G$\note{« $u : F \to G$ » en interligne.} de foncteurs $F, G : \mathcal{E} \to
(\mathrm{Ens})$ est dit une \emph{équivalence} s'il satisfait la condition
suivante :

($\alpha$) $\exists\, n \in \mathbf{N}$ tel que pour \struck{deux} tout
ensemble fini $E$ \struck{\ill{}} de cardinal $\geqslant n$,
\struck{et tt appl.\ \ill{} $E' \hookrightarrow E$ \ill{} $E$}, on ait
$u(E) : F(E) \xrightarrow{\sim} G(E)$.
\note{le passage biffé, sur une ligne et demie, ne se lit qu'en partie.}

On a

a) Si $F \xrightarrow{u} G \xrightarrow{v} H$ et si deux parmi les trois
$u, v, vu$ sont des équivalences, aussi la 3e.

b) L\struck{\ill{}}es équivalences sont stables par chgt de base

\page{26}
\[
  \underline{\Omega} = \underline{\mathrm{End}}(\mathcal{E}, \mathcal{E})
  \subset \underline{\mathrm{Hom}}(\mathcal{E}, \mathrm{Ens}) =
  \widehat{\mathcal{E}^{\circ}}
\]

$\underline{\Omega}$ est la sous-catégorie (pas pleine !) de
$\widehat{\mathcal{E}^{\circ}}$ formée :

a) comme objets, des foncteurs $F : \mathcal{E} \to (\mathrm{Ens})$ qui se
factorisent par la sous-catégorie $\mathcal{E}$ de $\mathrm{Ens}$, i.e.\ tels
que

\quad 1°) $F(E)$ est fini pour tt $E \in \mathrm{Ob}\,\mathcal{E}$

\quad 2°) $F(u)$ est injectif $\forall\, (u : E \to F) \in
\mathrm{Fl}(\mathcal{E})$

b) comme morphismes $F \xrightarrow{u} G$, des seuls homomorphismes tels que
$u(E) : F(E) \to G(E)$ soit injectif pour tt $E \in \mathrm{Ob}\,\mathcal{E}$.

$\underline{\Omega}$ a des opérations $+$, $\times$, $\circ$ :
\[
  \begin{cases}
    (F + G)(E) = F(E) \sqcup G(E) \\
    (F \times G)(E) = F(E) \times G(E) \\
    (F \circ G)(E) = F(G(E))
  \end{cases}
\]
\note{le feuillet s'arrête là ; le bas est vierge.}

\section{[Opérations $\lambda$ et représentations des groupes symétriques]}

\page{27}
Sur Topos loc.\ annelé (commutativement) $X$, on a
\[
  K(X, *) = \sum_{n \geqslant 0} K(B_{\mathfrak{S}_n X})
\]
On a pour $\xi \in K(X)$
\[
  \tau^*(\xi) = \sum \tau^n(\xi) \in 1 + K(X, *)^{\wedge +}
  = 1 + \xi + \tau^2\xi + \cdots
\]
Pour $\xi = \mathrm{cl}_{K(X)}(E)$, on a
$\tau^n(\xi) = \mathrm{cl}_{K(B_{\mathfrak{S}_n X})}
\bigl(\widetilde{\otimes}^{n} E\bigr)$.
\note{le signe de produit tensoriel porte un tilde et un exposant peu net ;
« $\widetilde{\otimes}^{n}$ » est une lecture proposée.}

Ceci posé, on a la relation
\[
  \tau^n(\xi) =
  \sum_{\substack{(\alpha_1, \ldots, \alpha_n) \in \mathbf{N}^n \\
  1\alpha_1 + 2\alpha_2 + \cdots + n\alpha_n = n}}
  \Sigma_{\alpha_1, \ldots, \alpha_n}(\lambda^1\xi, \ldots, \lambda^n\xi)\,
  \rho_1^{\alpha_1} \cdots \rho_n^{\alpha_n}
\]
\marginal{Si $\xi$ est effectif, s'agit-il d'une somme directe ? (Relation
avec le th.\ de commutation de H.\ Weyl ($\longrightarrow$ \ill{} ?))}
\note{note marginale écrite en oblique dans la marge gauche, en regard de la
formule.}
où $\rho_i \in K(B_{\mathfrak{S}_i}) = R_{\mathbf{Z}}(\mathfrak{S}_i)$ est
l'élément unité de $K(B_{\mathfrak{S}_i})$, regardé comme él.\ de degré $i$
dans
\[
  R(1, *) = K(e, *) = \sum_{i \geqslant 0} R(\mathfrak{S}_i)
\]
\note{sous $e$, une flèche et les mots « topos ponctuel ».}
(de sorte que $\rho_1^{\alpha_1} \cdots \rho_n^{\alpha_n}$ est un élément de
$R(1, n) = R(\mathfrak{S}_n)$), et \struck{$\Sigma_{\alpha_1, \ldots,
\alpha_n}($}
\[
  \Sigma_{\alpha_1, \ldots, \alpha_n} \in \mathbf{Z}[S_1, \ldots, S_n]
\]
est le polynôme isobare de poids $n$ ($S_i$ étant de poids $i$) qui
exprime la fonction symétrique à $n$ variables
\[
  \sum (x_1 \cdots x_{\alpha_1})(x_{\alpha_1 + 1} \cdots
  x_{\alpha_1 + \alpha_2})^2 \cdots
  (x_{\alpha_1 + \cdots + \alpha_{n-1} + 1} \cdots
  x_{\alpha_1 + \cdots + \alpha_n})^n
  = \sum x_1^{\nu_1} \cdots x_m^{\nu_m}
\]
(avec $\alpha_i$ = nb des $\nu_j$ égaux à $i$) comme polynôme en les fonctions
sym.\ élémentaires $S_1, S_2, \ldots, S_n$ en
\note{l'insertion « isobare de poids $n$ … » est de sa main, en interligne ; la
phrase s'interrompt au bas du feuillet sur « en », et le feuillet 28 ne la
reprend pas.}

\page{28}
\emph{Corollaire.} \struck{On a} $\lambda^n(\xi)$ apparaît comme coefficient
de $\rho_1^n$ dans $\tau^n$ ; $S^n(\xi)$ (opération de Adams) comme
coefficient de $\rho_n$.

[\emph{NB} Il est facile sans doute de vérifier que les
$1_{K(X)} \cdot \rho_1^{\alpha_1} \cdots \rho_n^{\alpha_n}$ sont
lin.\ indép.\ sur $K(X)$ dans $K(X, *)$ …]

\emph{Question.} Pour quels anneaux $k$ est-il vrai que
\[
  R_k(1, *) = \sum_{i \in \mathbf{Z}} R_k(\mathfrak{S}_i)
\]
est l'anneau de polynômes engendré par les $\rho_i$ ($i \geqslant 1$) ?
Est-ce vrai si $k$ est un anneau local (plus gén., tel que
$K(k) \xleftarrow{\sim} \mathbf{Z}$) ?
\note{l'indice de sommation sous $R_k(\mathfrak{S}_i)$ se lit
« $i \in \mathbf{Z}$ » ; on attendrait $i \geqslant 0$.}

\section{[Le jeu de go]}

\page{29}
jeu de go complet (sur go-ban $\nu \times \nu$)

noir joue $N$ coups

blanc joue $N'$ coups, \quad $N' = N$ ou $N - 1$
\note{le « $-$ » de « $N - 1$ » est surchargé ; il se lit aussi « $+$ ».}

\[
  \boxed{N = P + Q} \quad \boxed{N' = P' + Q'} \quad \text{i.e.} \quad
  \boxed{\varepsilon \overset{\text{déf}}{=} N - N' = 0 \text{ ou } 1}
\]
\note{sous le 0, une flèche : « si blanc joue dernier ».}

$P$ nb de pierres noires sur go-ban en fin de partie

$P'$ \quad ——— blanc ———

$Q$ nb de prisonniers noirs

$Q'$ nb de prisonniers blancs

$T$ territoire noir

$T'$ territoire blanc
\[
  \boxed{P + P' + T + T' = \nu^2}
\]
pts de noir \quad $Q' + T = S$

pts de blanc \quad $Q + T' = S'$
\note{devant $S$ et $S'$, une lettre noircie, biffée.}

différence
\begin{align*}
  D \ (= S - S') &= (Q' + T) - (Q + T') = (Q' - Q) + (T - T') \\
  &= (P - P') + (T - T') - \varepsilon
\end{align*}
\[
  D + \nu^2 = 2(P + T) - \varepsilon \qquad \text{ou} \qquad
  \nu^2 - D' = 2(P' + T') + \varepsilon \quad
  (D' \overset{\text{déf}}{=} -D = S' - S)
\]
\note{le signe devant $\varepsilon$, dans ces lignes, est un gros trait
noirci ; on le lit « $-$ », ce que le calcul demande. Dans la deuxième
formule, « $\nu^2 - D'$ » est bien ce qui est écrit ; c'est
$\nu^2 - D$ qui vaut $2(P' + T') + \varepsilon$.}
\[
  \boxed{D = 2(P + T) - \nu^2 - \varepsilon}
\]
\[
  (D' = 2(P' + T') - \nu^2 + \varepsilon)
\]

\page{30}
Supposons qu'il y a des territoires constitués de blanc et noirs considérés
\uncertain{inattaquables}, blanc a joué dernier (donc c'est à noir de
jouer) — .

pour $n$ coups
\[
  n = p + q \qquad n' = p' + q' \qquad n = n'
\]
$p$ nb des pierres noires sur goban

$q'$ \uncertain{nb} \ill{} prisonniers virtuels

$q$ nb de prisonniers virtuels

$t$ territoires noir (une fois enlevés les pris.\ bl.)

$t'$ \quad — \quad blanc ( ——— noirs)

$p + p' + t + t' = \nu^2$

$s = t + q'$ \quad pts de noir conventionnels

$s' = t' + q$ \quad ——— blanc ———
\[
  d = s - s' = (t - t') + (q' - q) = (t - t') + (p - p')
\]
\[
  \nu^2 + d = 2(p + t) \quad \text{i.e.} \quad
  \boxed{d = 2(p + t) - \nu^2}
\]
\[
  D = d - \varepsilon \qquad \varepsilon = 0, 1 \qquad
  (\text{NB } p + t = P + T)
\]

! (Comme on ne peut prévoir la valeur de $\varepsilon$ (est-elle indépendante
de la suite du jeu ?) il est raisonnable de compenser l'avantage du noir (du
commencement) en posant $\varepsilon = 1$ (i.e.\ en admettant que noir joue
dernier)).

\page{31}
Une pierre ne peut se placer en position de “prise”, i.e.\ sur une
intersection libre telle qu'il n'ait plus \struck{de} aucune liberté
(i.e.\ que le groupe qui la contient n'ait plus de liberté) sauf dans les deux
cas suivants :

1°) Cas du “Ko” cf.\ plus bas

2°) Cas de “remplissage d'un trou” : \struck{disons que} $A$ joue, sur une
\struck{case} intersection $i$ qui est la seule liberté restante d'un
ensemble $E$ réunion de groupes de $B$ (groupement de $B$), de
façon que toutes les libertés des groupes $G$ de $A$ contenant $i$ soient
prises par des pierres de $E$. ($E$ est entouré par les pierres de $A$ sauf
en $i$, et $E$ entoure le groupe $G$ de $A$ contenant $i$)
\note{en interligne, de sa main : « aucune » (au-dessus de « de », biffé),
« intersection $i$ » (au-dessus de « case »), « $E$ réunion » et « (groupement
de $B$) ». Plus bas, de même : « Les » au-dessus de « Chaque », « un des »,
« intersections $c$ » au-dessus de « cases », et le mot douteux « isolées ».}

\struck{\uncertain{Comprenant les}} \struck{Chaque} Les joueurs jouent
\emph{à tour de} rôle, et ils ne peuvent sauter leur tour (obligation de
jouer). L'obligation de jouer tombe dans le seul cas suivant :

(Fin de partie) Le joueur $A$ n'a plus de case libre où il soit légitime de se
mettre (i.e.\ il n'y soit pas en état de prise, sauf l'exception un des
1° et 2° plus haut), sauf des \struck{cases} intersections $c$ du type
suivant : il y a un groupement $E$ de $A$ n'ayant plus que deux libertés
\struck{intérieures} \uncertain{isolées} \struck{\ill{}}

\page{32}
\struck{entourées de pierres de $B$.}
(toutes les autres libertés étant prises par des pierres de $B$)
\struck{et alors \ill{} serait nécessairement entouré de pierres de $A$}
(et $c$ est une de ces deux-là).
[\emph{NB} Si $A$ s'y mettait, dans un prochain coup $B$ en
\struck{\ill{}} plaçant dans la dernière liberté $c'$ disponible du groupement
$E' = E \cup \{c\}$, capturerait $E'$ — et de proche en proche il capturerait
toutes les pierres de $A$ !], même si $A$ avait un avantage considérable sur
$B$ par les territoires conquis !]
On exige de plus, cependant, que le groupement $E$ soit “inattaquable”, i.e.\
que si c'était le tour de $B$ de jouer au lieu de $A$, il ne pourrait se
placer ni en $c$, ni en $c'$ (i.e.\ que l'on ne peut trouver un
sous-groupement $E_c$ de $E$ entourant $c$ \struck{dont} $c$ soit la
seule liberté \struck{\uncertain{intérieure}} restante,\note{« entourant $c$ » et la parenthèse qui suit sont en interligne.} (ni un
sous-groupement $E_{c'}$ de $E$ entourant $c'$ dont $c'$ soit la seule liberté
restante), ni que $B$ puisse jouer en $c'$ par la règle du Ko.

\page{33}
En fin de partie

noir a $\rho$ groupements solidaires (où chacun $G_i$ a $y_i \geqslant 2$
yeux)

blanc : $\rho'$ groupements solidaires ( … $y'_j \geqslant 2$)

\struck{① si blanc joue dernier ($\varepsilon = 0$) alors les $\rho'$
groupements ont chacun exactement 2 yeux, donc $T' = 2\rho'$}
\note{passage encadré et barré de traits obliques.}

si $t - 2\rho \leqslant t' - 2\rho'$

alors c'est blanc qui joue le dernier (i.e.\ $\varepsilon = 0$, i.e.\
$D = d$)

si $t' - 2\rho' < t - 2\rho$,

c'est noir qui joue le dernier (i.e.\ $\varepsilon = 1$, i.e.\ $D = d - 1$)
\note{le texte occupe la moitié droite d'un feuillet double ; la moitié gauche
est vierge.}

\section{Groupes formels}
\note{les mots du titre sont de sa main, seuls, en haut à droite du feuillet
34, qui ne reçoit pas de numéro de page ; le feuillet 35 les répète en tête.}

\page{35}
\emph{Groupes formels}

(I) Extension d'un $p$-divisible par un lisse (ou unipotent) est triviale

\emph{Cor.} Extension d'un $p$-divisible a-torique par un lisse (ou
unipotent) est \struck{\ill{}} \uncertain{un produit direct}
\note{au-dessus de la fin de la ligne, trois mots dont les derniers sont
biffés et surchargés à l'encre plus foncée.}

[Tout sous-groupe $p$-divisible a-torique est central.

Tout groupe quotient \struck{a-toriques} $p$-div.\ l'est
\struck{\ill{}} \ill{} avec le plus grand sous-groupe $p$-divisible
a-torique — \uncertain{quand} \uncertain{ce dernier} est \uncertain{central}
…

\ill{} \ill{} sont des \uncertain{plus grands} sous-groupes
$p$-divisibles a-toriques :
\[
  \left\{\begin{array}{l}
    \text{existe} \\
    \text{stabilité par ext.\ corps de base} \\
    \text{indépendance vis-à-vis sous-groupes} \\
    \text{sur groupes quotients}
  \end{array}\right.
\]

\emph{NB} Pour certains groupes formels sur \struck{a-}$p$-div.\
toriques il \uncertain{faut}).]
\note{en interligne, de sa main : « $p$-div. » (deux fois, au-dessus de
« a-toriques » et de « a- », biffés), « sur groupes quotients » ; plus bas,
« a-torique », « on a » et « isogène ».}
\marginal{il suffit \uncertain{aussi} de le faire pour
$(G/R)_{0s}$ \uncertain{et $G/Z$} \uncertain{linéaire}. \emph{Mieux} :
$\mathrm{Cent}\,(G/\mathrm{Cent}\,G)$ est quasi-unipotent}
\note{note écrite verticalement dans la marge gauche, montant le long du
crochet ; une flèche la relie à la ligne « Rem. » qui suit. Lecture très
incertaine au-delà des formules.}

\emph{\uncertain{Rem.}} $(G_{0s})/I_{\ill{}}\ \ill{} = \pi_0$ \ill{} par des
produits a-toriques i.e.\ est \emph{linéaire}

(II) Extension d'un groupe formel [lisse] par un $p$-divisible est
[i.e.\ $G$ conn.] \struck{triviale} un produit \ill{} — à isogénie près des
\uncertain{\ill{}} …

\emph{Cor.} $DG \cap \mathfrak{Z}(G)$ est
\struck{\ill{}} quasi-unip., \uncertain{i.e.}\ primaire au $p$.
[\emph{Mieux}, \struck{\ill{}} $\mathrm{Cent}(\mathrm{Dér}\,G)$ est
quasi-unip.]
\struck{\ill{} $\mathrm{Dér}\,G$, \ill{} $\mathfrak{Z}(G)$}, \ill{}

[\emph{Dém.} Soit \struck{\ill{}} $T$ le plus grand
\uncertain{sous-groupe} $p$-divisible a-torique de $G$, on a
$G \simeq T \times G/T$, donc $DG \cap T = \{e\}$]

\emph{Cor.} $G$ (lisse) est \uncertain{isogène} \struck{\ill{}} à
\uncertain{un produit} de la \ill{} produits a-toriques par un facteur … ?

\emph{Question.} $DG \cap \mathfrak{Z}(G)$ est-il $\simeq$
unipotent, i.e.\ ne contient pas de $\mu_p$ (\uncertain{base} \ill{}) ?

\emph{Solution}

\begin{tikzcd}[column sep=small, row sep=small]
 & L_1 \arrow[r, "d"] & L_0 \arrow[dr] & \\
1 \arrow[r] & \Pi_1 \arrow[u] \arrow[d] & & \Pi_0 \arrow[r] & 1 \\
 & D_1 \arrow[r, "\int"'] & D_0 \arrow[ur] &
\end{tikzcd}

\note{le schéma est redessiné : $1 \to \Pi_1$ se sépare en deux flèches
obliques vers $L_1$ et $D_1$, qui convergent de même, à droite, de $L_0$ et
$D_0$ vers $\Pi_0 \to 1$ ; sous $D_0$, entre crochets, une flèche
$\downarrow \theta$ vers $\mathrm{Aut}\,D_1$.}

\[
  \begin{bmatrix}
    L_1 = \mathfrak{Z}(G) \\
    D_1 = D(G) \\
    L_0 = G/\mathfrak{Z}(G) \\
    D_0 = G/D(G)
  \end{bmatrix}
  \qquad
  \begin{cases}
    D_1^{D_0} = \Pi_1 \\
    D_{1\,D_0} = \{e\}
  \end{cases}
  \qquad
  \begin{array}{l}
    \Pi_1 = L_1 \cap D_1 \\
    \Pi_0 = G/L_1 D_1
  \end{array}
\]

$p^{\infty}\Pi_1 = e$

$\mathrm{at}(\Pi_0) = \{e\}$

+ extension \uncertain{\ill{}} à \uncertain{quasi-iso} \ill{}
\[
  \underbrace{\Pi_1, \ \mathrm{Ind} \times \mathrm{Ind}\,\delta, \ \Pi_0}
  \qquad
  D' = D_1 \amalg_{\Pi_1} L_1 \qquad L' = L_0 \times_{\Pi_0} D_0
\]
\marginal{\emph{NB} $\mathrm{Im}\,\delta \simeq \hat{H}$, $H$ grp alg.\
\uncertain{canonique}, \ill{}}
[obstruction nulle dans $H^3(\Pi_0, \Pi_1)$

indétermination dans $H^2(\Pi_0, \Pi_1)$]

Traiter cas $\Pi_0$ $p$-divisible (plus gén.\ des extensions des sous-groupes
de $G$ images inverses de $p^{\infty}\Pi_0$

\emph{Dire quels sont les} $L_1 \xrightarrow{d} L_0$ possibles
$(D_1 \xrightarrow{\int} D_0, \theta)$ —
\marginal{(conditions \uncertain{sur} $\Pi_1$, $\Pi_0$ \struck{\ill{}})}

\emph{Conditions pour que $G$ soit} \uncertain{linéaire} $DG$
soit \uncertain{linéaire} \uncertain{can.}\ + algébrisable …
\note{au bas à droite, un griffonnage de formules surchargées où se lisent
« $= \Pi_1$ » et « $D_0 = \{e\}$ » ; en bas à gauche, « TSVP ». Le feuillet 36
fait suite.}

\page{36}
Dans la situation générale des ext.\ par $G$ \uncertain{correspond} à une
situation
\[
  1 \to \Pi_1 \to G_1 \xrightarrow{d} G_0 \to \Pi_0 \to 1, \qquad
  G_0 \overset{\theta}{\dashrightarrow} \mathrm{Aut}(G_1, \Pi_1)
\]
on a (pour \uncertain{la} $G$ donnée) un hom
\[
  (\mathrm{Cent}\,G_0) \cap \mathrm{Ker}\,\theta \xrightarrow{\varphi_G}
  Z^1(\Pi_0, \Pi_1)
\]
dont le noyau est l'image de $Z(G)$ dans $G_0$.
\marginal{$G \leftrightsquigarrow (\Pi_0, \mathrm{Ind}, \Pi_1)$, $G^{\bullet}$}
\note{dans la marge de droite, deux accolades obliques relient $G$ et
$G^{\bullet}$ au triplet ; la flèche est de l'éditeur.}

Dans le cas qui nous intéresse
\[
  [\,G_1 = D' = L_1 . D_1 = \mathrm{Cent}\,G . \mathrm{Dér}\,G, \quad
  G_0 = L' = G / L_1 \cap D_1 = G/\mathrm{Cent}\,G \cap \mathrm{Dér}\,G
  \ \text{--}\,]
\]
cela \struck{\ill{}} s'exprime par
\[
  \mathrm{Cent}(D_0) \cap \mathrm{Ker}\,\theta \xrightarrow{\varphi_G}
  \mathrm{Hom}(\Pi_0, \Pi_1)
\]
(\uncertain{ce qui exprime} que $L_0 \hookrightarrow \mathrm{Cent}\,G$ est
\uncertain{iso}), dont le noyau est nul ; ceci montre donc qu'on doit avoir
(compte tenu que $\Pi_1$ est quasi-unipot.)
\[
  \boxed{\mathrm{Cent}\,D_0 \cap \mathrm{Ker}\,\theta \text{ est
  quasi-unipotent !}}
\]
\struck{$\sqrt[2]{\,[\mathrm{Cent}(\mathrm{Dér}(G)) / \mathrm{Cent}(G)]}
\cap \mathrm{Cent}\,D_0$}

En fait, on doit avoir \uncertain{$G$ unipotent}, comme
$D_0 \simeq G/(G, G)$, $\mathrm{Cent}(D_0)$ \emph{unipotent}.

\begin{tikzcd}[column sep=small, row sep=small]
 & & L_1 \arrow[dr, "\mathcal{H}"] & & \\
1 \arrow[r] & \Pi_1 = L_1 \cap D_1 \arrow[ur, hook, "\mathcal{A}"] \arrow[dr, hook', "\mathcal{H}"'] \arrow[rr, "\mathcal{A} * \mathcal{H}"] & & D' \arrow[r, hook] & G \\
 & & D_1 \arrow[ur, "\mathcal{A}"'] & &
\end{tikzcd}

\note{les étiquettes $\mathcal{A}$ et $\mathcal{H}$ des flèches sont
accompagnées d'un trait ondulé ; un arc relie $L_1$ à $D_0$, un autre $D_1$ à
$L_0$, au-delà de $D'$. À gauche de $1$, un mot biffé.}

\page{37}
\note{Feuillet de brouillon ; on en donne les éléments dans l'ordre de la
page.}

$G/\mathbf{Z}$ \qquad
$p^{\infty}\,\mathrm{Cent}(G/\mathbf{Z})$ \qquad
$\underline{\mathrm{Hom}}_{\mathrm{gr}}(T, \mathbf{Z}) \uparrow G$ \qquad
\struck{\ill{}} $\underbrace{\mathbf{Z} \times' T}$

$\mathrm{Cent}(G/\mathbf{Z})$ \qquad \struck{$\mathrm{Cent}(G/\mathrm{Cent}$}

\[
  \left\{
  \begin{array}{l}
    \boxed{\mathrm{Cent}(\mathrm{Dér}\,G) \text{ unipotent}} \\
    \boxed{\struck{\ill{}}\ (G/\mathbf{Z})_{\mathrm{ab}}
      \ \struck{\text{a-tori}\ill{}} \ \uncertain{\text{linéaire}}} \\
    \boxed{\mathrm{Cent}(G/\mathbf{Z}) \text{ unipotent}}
  \end{array}
  \right.
\]
\[
  (\mathrm{Dér}\,G)_{\mathrm{ab}} \text{ unipotent ?} \longrightarrow
  \text{voir fin}
\]
\struck{$\mathrm{Dér}\,G_{\mathrm{ab}} = \ill{}$}
\note{la dernière question est entourée ; la flèche vers « voir fin » en part.
Le $\mathbf{Z}$ est ici le gros Z au trait double, peut-être le centre.}

\page{38}
\note{Feuillet de brouillon ; une ligne coupée en haut du feuillet ne se lit
pas.}

\begin{tikzcd}[column sep=small, row sep=small]
 & \mathfrak{Z} \arrow[dr, no head] & & & \\
Z \cap DG \arrow[ur, no head] \arrow[dr, no head] & & \mathfrak{Z}.DG \arrow[r, no head] & G \arrow[r, no head] & \mathcal{G} \\
 & DG \arrow[ur, no head] & & &
\end{tikzcd}

\note{$H$ est écrit à gauche de $DG$ ; les traits du losange sont
sans pointe.}

\[
  \underbrace{Z \cap DG \to DG \to Z.DG}
  \to G
\]

$DG$ \qquad $G/\mathbf{Z}$ \qquad
$DG / Z \cap DG$

\[
  1 \to \mathfrak{z} \to D \to \hat{H} \to 1
\]
$H$ groupe alg.\ linéaire connexe (lisse)

$\mathfrak{z}$ groupe formel \struck{\ill{}} (lisse)

$D$ extension centrale de $\hat{H}$ par $\mathfrak{z}$
\struck{\ill{}}
\marginal{groupe formel \uncertain{commutatif}}

$\mathfrak{z} \to \mathfrak{Z}$ hom.\ de groupes formels
\uncertain{commutatifs} \ill{}
\note{la note marginale est rattachée à $\hat{H}$ par une accolade ; une
flèche part du passage biffé vers elle.}

\page{39}
\note{Feuillet de brouillon à l'italienne, couvert de schémas ; on en donne
les éléments dans l'ordre de la page, sans les relier.}

\begin{tikzcd}[column sep=small, row sep=small]
H \arrow[d, hook'] & H_T \arrow[l] \arrow[d, hook'] \\
G \arrow[d] & G_T \arrow[l] \arrow[d] \\
I & T \arrow[l]
\end{tikzcd}

\note{à droite de ce carré : $H_T^i$, $W \supset W_i$, $W$ ; sous $T$ :
$\check{T}_{a_i}$ ; un trait l'enferme et le relie à la suite
$H \to E \to G \to 1$.}

$H \to E \to G \to 1$ \qquad $\varphi(x, y)$ \qquad
$\psi(x + y) - \psi(x) - \psi(y) = \varphi(x, y)$

\struck{$\varphi(x, y)$} $(\varphi(x, y))^{-1} = \varphi(x, y)$
\note{lecture très incertaine.}

$1 \to \hat{G}_a \to E \to \hat{G}_a \to 1$ \qquad
$Z(E) = D(E)$

\[
  \begin{array}{c} [Z \to A] \\ \wr\!\wr \\ {[D \to L]} \end{array}
  \qquad
  \begin{array}{c} [\mathfrak{Z} \to A] \\ \wr \\ {[\hat{D} \to \hat{L}]}
  \end{array}
  \qquad
  \begin{array}{c} \mathfrak{Z} \to A \\ \wr\ \ \ \wr \\ \hat{D} \to \hat{L}
  \\ \mathfrak{Z} \to A \\ D \to L \end{array}
\]
\note{sous les deux premiers crochets, un chapeau commun ; plus bas, deux
schémas en losange : $\Pi_1 \to (Z \to A,\ D \to L) \to \Pi_0$ et
$\Pi_1 \to (\mathfrak{Z} \to A,\ \hat{D} \to \hat{L}) \to \Pi_0$.}

\struck{$\Phi(x, y) =$}
\[
  \Phi(x, y) = \underbrace{\varphi(x, y)}_{\ill{}} +
  \underbrace{\psi(x + y) - \psi(x) - \psi(y)}
\]
$H^2(\hat{G}_a, \hat{G}_a)$

\end{document}
